%! Modeling with Exponential Functions: Interest, Half-Life, and e — Unit 6, Lesson 6.4 — Student Packet Cover Sheet
\documentclass[10pt]{article}
\usepackage{algebra2-article}
\usepackage{algebra2-boxes}

\begin{document}

% Full-bleed forest banner
\begin{tikzpicture}[remember picture, overlay]
  \fill[forest] (current page.north west)
    rectangle ([yshift=-0.9in]current page.north east);
\end{tikzpicture}
\vspace{-0.8in}

\begin{center}
  {\color{white}\textbf{\LARGE Algebra 2}}\\[4pt]
  {\color{forestmid}\large\textbf{Unit 6 \quad Exponential Functions}}\\[6pt]
  {\color{white}\textbf{\large Lesson 6.4 \quad Modeling: Compound Interest, Half-Life, and $e$}}
\end{center}

\vspace{0.22in}
\namedateperiod
\vspace{0.18in}

\begin{learningtargetbox}
\small
\begin{enumerate}[leftmargin=1.5em, itemsep=5pt]
  \item \ldots{} build and evaluate an exponential model for \textbf{compound interest}, \textbf{half-life}, and \textbf{doubling time} --- and read the answer back as a sentence about money, medicine, or population.
  \item \ldots{} explain why the base is always the factor for \textbf{one period} while the exponent \textbf{counts periods}, so the rate gets divided by $n$ and the time gets multiplied by it.
  \item \ldots{} describe where the number $e$ comes from and use $A=Pe^{rt}$ for continuous compounding.
  \item \ldots{} use the \textbf{range} and the \textbf{horizontal asymptote} to say what a decay model never does, and recognize an equation whose answer exists but that no common base will ever solve.
\end{enumerate}
\end{learningtargetbox}

\vspace{0.14in}

\begin{tocbox}
\small
\renewcommand{\arraystretch}{1.5}
\begin{tabularx}{\linewidth}{@{} c l X r @{}}
\rowcolor{forestbg}
\textbf{\#} & \textbf{Component} & \textbf{Description} & \textbf{Score} \\
\midrule
1 & Warm-Up        & Spiral review: splitting a yearly rate, counting periods, halving, and one equation with no common base & \blank{1.2cm} \\
2 & Guided Notes   & Compound interest, the number $e$, half-life and doubling time, and the doubling wall & \blank{1.2cm} \\
3 & Group Activity & \emph{The Base Is One Period} --- setup tables, two offers compared, error analysis, and the Rule of 72, by tier & \blank{1.2cm} \\
4 & Exit Ticket    & Quick check: one interest model, one half-life model, an SOL-style item, and a question about the wall & \blank{1.2cm} \\
5 & Homework       & Four accounts, quarterly vs.\ continuous, caffeine, a doubling read off a graph, error analysis, and a depreciating truck & \blank{1.2cm} \\
\midrule
  & \multicolumn{2}{r}{\textbf{Total}} & \blank{1.2cm} \\
\end{tabularx}
\end{tocbox}

\vspace{0.10in}

\begin{remindbox}
\small Everything today is Lesson 6.1's $y=ab^{x}$ in a costume, and one sentence keeps the costumes
straight: \textbf{the base is what happens in one period, and the exponent counts the periods.} For
compound interest, $A=P\left(1+\frac{r}{n}\right)^{nt}$ --- the annual rate is \emph{divided} by $n$
because one period only earns a fraction of it, and the time is \emph{multiplied} by $n$ because there
are that many periods. Break the sentence in half and you get the two errors this lesson is made of:
using $1.08$ with an exponent of $20$ quarters (a full year's interest every quarter) or $1.02$ with an
exponent of $5$ years (five quarters of growth). Chopping the year finer and finer does not run away;
it climbs toward a ceiling, and that ceiling is the number \textbf{$e\approx2.71828$}, which turns
$A=P\left(1+\frac{r}{n}\right)^{nt}$ into $A=Pe^{rt}$. Half-life and doubling time are the same
structure on a different clock: $A_{0}\left(\frac12\right)^{t/h}$ and $A_{0}2^{\,t/d}$, where dividing
by $h$ or $d$ converts hours or years into \emph{periods} --- which is why a fractional exponent like
$t/h=1.5$ is normal and not a mistake. Two Lesson 6.0 facts stay in force: a decay model has range
$(0,A_{0}]$ with a horizontal asymptote at $0$, so half of something is never nothing. Finally, the
lesson ends on a question it cannot answer. When does \$1000 at $5\%$ double? The graph says between
year 14 and year 15, so the answer \textbf{exists} --- but $(1.05)^{t}=2$ has no common base, and $2$
is not a power of $1.05$. That is not ``no solution.'' It is the gap Unit 7 was invented to close.
\end{remindbox}

\end{document}
